\documentclass[11pt]{article}
\usepackage[a4paper,margin=26mm]{geometry}
\usepackage{amsmath,amssymb,amsthm}
\usepackage[T1]{fontenc}
\usepackage{lmodern}
\usepackage{hyperref}
\newtheorem{proposition}{Proposition}
\title{Noninteger roots of a Schubert counting polynomial\\\large Counterexample to TLMC 00000008430}
\author{ziangni-sys}
\date{4 October 2026}
\begin{document}
\maketitle
\begin{abstract}
The full Grassmannian of lines in a three-dimensional vector space over
$\mathbb F_q$ has $q^2+q+1$ points. This is a Schubert configuration
family determined by dimensions. Its counting polynomial is monic but
has nonreal roots, disproving the assertion that its roots are integers.
\end{abstract}

\section{Configuration family and scope}
The source defines a Schubert configuration as a chain of subspaces
determined by dimension conditions. It asserts that its finite-field
count is a monic polynomial in $q$ whose roots are integers.

Take $V=\mathbb F_q^3$ and the family of chains
\[
 0\subsetneq L\subsetneq V,\qquad \dim L=1.
\]
Each chain is uniquely determined by the line $L$, so this family is
the Grassmannian $\operatorname{Gr}(1,3)=\mathbb P^2$. It is the full
Schubert variety, namely the closure of the largest Schubert cell for
any fixed complete flag. The entire Grassmannian is allowed by the
dimension-chain definition in the source. We count this full variety,
not only its open affine cell, which by itself would have $q^2$ points.

\section{Count and polynomial}
\begin{proposition}
The number of these configurations over every finite field of cardinality
$q$ is $q^2+q+1$.
\end{proposition}
\begin{proof}
Every nonzero vector of $V$ spans a unique line. Each line has exactly
$q-1$ nonzero vectors, and the nonzero-vector sets of distinct lines
are disjoint. Thus the number of lines is
\[
 \frac{q^3-1}{q-1}=q^2+q+1.
\]
This holds for every finite field, not merely for selected small values.
The polynomial $P(X)=X^2+X+1$ is also uniquely determined by these
counts: any second counting polynomial agrees with it on every prime
field cardinality, and there are infinitely many primes. Two distinct
polynomials cannot agree at infinitely many points.
\end{proof}

\newpage
\begin{proposition}
$P$ is monic and has a root that is not an integer.
\end{proposition}
\begin{proof}
Its leading coefficient is $1$. Put
\[
 \zeta=\frac{-1+i\sqrt3}{2}.
\]
Direct expansion, using $(\sqrt3)^2=3$, gives
$\zeta^2+\zeta+1=0$. Since
$\operatorname{Im}\zeta=\sqrt3/2>0$, the root $\zeta$ is not a real
number and in particular is not an integer.
\end{proof}

The polynomial-counting and monicity clauses both hold in this example;
it is precisely the integer-root clause that fails. Thus the example
disproves the conjecture without imposing any alternate counting model.

\section{Lean correspondence and verification}
\raggedright
The Lean4.19.0 project pins Mathlib commit
\texttt{c44e0c8ee63ca166450922a373c7409c5d26b00b}.
For any field $k$, the type \texttt{Lines k} consists of actual
\texttt{Submodule k (Fin 3 -> k)} objects whose \texttt{Module.finrank}
is one. Fixing the endpoints $0$ and $k^3$ identifies these objects
exactly with the chains above; no extra choices or quotients change
the count.

The theorem \texttt{count\_lines} applies Mathlib's proved equivalence
between projective points and one-dimensional submodules, followed by
its projective-space cardinality theorem and the actual dimension of
$k^3$. It proves the count for every finite field. The theorem
\texttt{polynomial\_counts} identifies these cardinalities with
evaluation of the actual complex polynomial \texttt{countPolynomial}.
The theorem \texttt{counting\_polynomial\_unique} uses actual prime
fields \texttt{ZMod p}, infinitude of primes, and the polynomial
identity theorem to prove uniqueness of the counting polynomial.

The theorem \texttt{countPolynomial\_monic} checks monicity.
The value \texttt{zeta} is an actual complex number with the real square
root of three in its imaginary part. The theorems
\texttt{zeta\_isRoot} and \texttt{zeta\_not\_integer} prove its root
equation and exclusion from every complex cast of an integer.
The final theorem \texttt{conjecture\_00000008430} combines the
general counting identity with existence of a noninteger root.

Run \texttt{lake build} and \texttt{lake env lean Main.lean} inside
the \texttt{lean} directory. The axiom audits for the count, uniqueness,
monicity, and final theorem use only the standard logical axioms
\texttt{propext}, \texttt{Classical.choice}, and \texttt{Quot.sound}.
No table of counts, unproved geometric assertion, native evaluation,
numerical approximation, or auxiliary computational script is used.
\end{document}
